\section{Conclusions and staged research outlook}
\label{sec:conclusions}

The progression from classical curvature-squared gravity to QQG makes renormalization and the independent couplings central to a broad theoretical space.  Its scalar, compact-object, and cosmological sectors answer different questions.  A useful first contribution can be made within the PS theory without simultaneously solving the physical-state problem of the full spin-two sector.  Its value depends on a precise definition of the measure, renormalization prescription, observables, and limits being tested.

This research design specifies a candidate periodic regulator, its HMC force, and free-field and step-scaling predictions.  The finite-lattice bound in appendix~\ref{app:coercivity} establishes integrability for that fixed-dimensional measure.  It does not establish continuum existence, ultraviolet completion of the full metric theory, or a Lorentzian probability interpretation.  The confirmatory program is designed to distinguish controlled scaling, a regulator or tuning failure, and an inconclusive range of accessible volumes.  No interacting production ensembles or cosmological likelihood fits are reported in this draft.

The black-hole discussion then makes the role of Euclidean methods precise: a regular thermal horizon can encode physical information without resolving a Lorentzian interior singularity.  The no-boundary discussion gives a complementary criterion: probabilities for physical histories can remain predictive where a classical extrapolation fails, provided the state, constraints and probability prescription are defined.  An arrow of time requires an additional analysis of fluctuations and coarse graining.  The cosmological discussion identifies distance--growth consistency and specific early-universe predictions as possible later tests, each tied to an explicit action and state prescription.  The final AI-assisted experiment targets calibrated distinguishability of dark-energy and modified-gravity controls, with reference calculations and measured computational cost.  Its connection to QQG requires the same action-to-observable matching as any other inference.  These connections motivate further calculation while preserving the independent completion criteria of the scalar stage.  The next deliverable is a validated scalar implementation and a pilot error-and-cost assessment; only then should production resources and stronger physical claims be committed.
